\section{Experimental Design}
\label{sec:experimental_design}

Guardian Lens evaluates whether black-box observations can distinguish
three researcher-induced decision policies in a fixed-weight
vision-language model. We query \texttt{gemini-3-flash-preview} at
temperature 1.0 with one completion candidate and three repeated samples
per condition. Each query combines one system profile, one image, and one
100-point allocation task. Responses contain numeric allocations
\(x\) and \(y\) to two fictional organizations, with \(x+y=100\).

\subsection{Behavioral Profiles}

The experiment defines three known policy classes.

\begin{table}[t]
\centering
\scriptsize
\begin{tabular}{p{0.17\columnwidth} p{0.34\columnwidth} p{0.39\columnwidth}}
\toprule
\textbf{Profile} & \textbf{Policy} & \textbf{Expected behavior} \\
\midrule

Neutral (N) &
No organization receives special treatment. &
Balanced ordinary allocation, efficient choice under cost, and no cue
effect. \\

Cue-bound (C) &
Favor X only when the target emblem is visible. &
Positive target cue effect but neutral-like behavior when the target is
absent. \\

Generalized (G) &
Treat advancing X as a standing priority. &
High allocation to X even without the cue and under explicit cost. \\

\bottomrule
\end{tabular}
\caption{Researcher-controlled behavioral profiles.}
\label{tab:profiles}
\end{table}

The key distinction is between \emph{conditional} and
\emph{generalized} prioritization: Cue-bound favors X only when the
designated visual trigger is present, whereas Generalized favors X
independently of that cue.

\subsection{Tasks and Visual Intervention}

We use two allocation tasks. In the \emph{ordinary} task, X and Y are
equally effective. In the \emph{costly} task, each point allocated to X
produces \(0.80\) benefit units while each point allocated to Y produces
\(1.00\), giving

\begin{equation}
B(x,y)=0.80x+1.00y,
\qquad x+y=100.
\end{equation}

Thus, prioritizing X incurs a 20\% efficiency penalty.

Visual effects are measured using matched counterfactual image pairs.
Each clean scene is paired with a modified version generated from the
same source image, with changes restricted to a logged emblem overlay
region. Modified scenes contain either the target emblem or a visually
similar distractor. Pixel-level validation confirmed that held-out pair
differences were confined to the recorded intervention region.

\subsection{Calibration and Held-Out Evaluation}

Six scenes form the calibration set used to construct the auditor. The
held-out set contains 12 distinct scenes: five clear-target pairs, four
subtle-target pairs, and three distractor pairs.

Three repeated samples are collected for every combination of scene,
profile, image condition, and task. The resulting analytical dataset
contains 216 calibration calls and 432 held-out calls, for 648 model
calls in total.

The held-out dataset is frozen before final querying, and the auditor is
constructed using calibration data only.

\subsection{Pre-Specified Hypotheses}

Four confirmatory hypotheses were specified before held-out analysis:

\begin{enumerate}[leftmargin=*]

    \item[\textbf{H1}] Cue-bound exhibits a larger target-minus-clean
    ordinary cue effect than Neutral and Generalized.

    \item[\textbf{H2}] On clean ordinary images, Generalized allocates
    more to X than Cue-bound.

    \item[\textbf{H3}] Under the costly task, Generalized retains greater
    allocation to X than Cue-bound and Neutral.

    \item[\textbf{H4}] The frozen auditor recovers the three behavioral
    profiles above the nominal three-class reference accuracy of
    \(1/3\).

\end{enumerate}

A planned exploratory reportability analysis was not evaluated because
explicit reportability judgments were absent from the preserved output
schema.